\section*{Chapter \ref{ch:mx:metaorders-latent-liquidity-and-cross-impact} --- Metaorders, Latent Liquidity and Cross-Impact}

\begin{solution}{exo:mx:metaorders-latent-liquidity-and-cross-impact:1}
$\sqrt{2\times50/2}=\sqrt{50}=7.07$ ticks.
\end{solution}

\begin{solution}{exo:mx:metaorders-latent-liquidity-and-cross-impact:2}
The latent orders diffuse back toward the price faster than the metaorder removes them, so it always meets a nearly full book: the price moves in
proportion to what it has taken, like a constant-depth book.
\end{solution}

\begin{solution}{exo:mx:metaorders-latent-liquidity-and-cross-impact:3}
With $I(t)=I_T\sqrt{t/T}$ and constant rate, the average is $\frac1T\int_0^TI_T\sqrt{t/T}\,dt=\tfrac23I_T$.
\end{solution}

\begin{solution}{exo:mx:metaorders-latent-liquidity-and-cross-impact:4}
It is not symmetric (0.5 against $-0.1$) and has a negative diagonal entry, so a trader could cycle through the two assets at a profit. The projection
symmetrises it and clips its negative eigenvalues to zero, giving the nearest symmetric positive semi-definite matrix.
\end{solution}

\begin{solution}{exo:mx:metaorders-latent-liquidity-and-cross-impact:5}
Selling both, each sale's \omterm{def:mx:metaorders-latent-liquidity-and-cross-impact:cross}{cross-impact} lowers the other's price further, a cost the own-impact model leaves out (36.0 against 52.1). In the pair trade,
selling the first lowers the second's price while it is being bought, a gain the model leaves out (36.0 against 20.0).
\end{solution}

\begin{solution}{exo:mx:metaorders-latent-liquidity-and-cross-impact:6}
The cost matrix is then a Kronecker product $\Lambda\otimes E$, whose inverse is $\Lambda^{-1}\otimes E^{-1}$: the optimum is $(\Lambda^{-1}\mu)\otimes
E^{-1}\mathbf 1$, the same time shape $E^{-1}\mathbf 1$ for every asset, scaled to each asset's total. Only the cost depends on the cross terms.
\end{solution}

\begin{solution}{exo:mx:metaorders-latent-liquidity-and-cross-impact:7}
With a 100-second execution the exponent stays at 1.0 up to $Q=100$ and falls to 0.94 and 0.75 only for 300 and 1\,000: the crossover moves up tenfold, to
about $LDT=100$, because the book has ten times longer to refill.
\end{solution}

\begin{solution}{exo:mx:metaorders-latent-liquidity-and-cross-impact:8}
The price falling below the average fill after the end is what fair pricing with a transient impact predicts, not proof of anyone's loss; the price after
the end keeps moving (in the latent book it keeps falling). Who lost depends on how long the fillers held and where they closed; measure their mark-outs
at their own horizons.
\end{solution}

\begin{solution}{pb:mx:metaorders-latent-liquidity-and-cross-impact:1}
\textbf{1.} Splitting with decaying impact, information and fair pricing, and a latent book that vanishes at the price.
\textbf{2.} Intentions to trade not placed in the book; a latent density $L|x-p|$.
\textbf{3.} $\int_0^ILx\,dx=LI^2/2=Q$.
\textbf{4.} 0.05, 0.18, 0.52, 1.74, 4.9, 12.2, 23.3 and 44.0 ticks; exponent 1.0 up to $Q=10$, then 0.94, 0.76, 0.59, 0.53: the crossover near $LDT=10$.
\textbf{5.} The average price paid equals the price after completion.
\textbf{6.} Two thirds of the peak.
\textbf{7.} 1.9 seconds after the end for $Q=10$ (2.3 for 300); it keeps falling, to 0.16 (0.27) of the peak after 100 seconds.
\textbf{8.} The latent orders never lose memory in this version: the book refills around the old price; deposition and cancellation of latent orders
would anchor part of the move.
\textbf{9.} $r_t=\Lambda q_t+\text{noise}$.
\textbf{10.} Imbalance in one stock or industry moves others' returns, persistently, often negatively; trades mediate much of the covariance of returns,
with liquidity sharing the correlations' sectorial structure.
\textbf{11.} $\begin{pmatrix}0.96&0.42\\0.42&0.81\end{pmatrix}$; 4.5\% and 5.3\% of return variance beyond each stock's own flow.
\textbf{12.} Linear, odd and symmetric \omterm{def:mx:metaorders-latent-liquidity-and-cross-impact:cross}{cross-impact}, for bounded kernels, or cycles through the assets would pay.
\textbf{13.} 52.1 against 36.0.
\textbf{14.} 20.0 against 36.0.
\textbf{15.} 64.6 for the sale and 38.4 for the pair: the joint schedule saves 19\% and 48\%.
\textbf{16.} When \omterm{def:mx:metaorders-latent-liquidity-and-cross-impact:cross}{cross-impact} decays at a different rate from own impact; with a slower cross decay the pair trade saves a further 3.7\%.
\textbf{17.} \emph{Named result}: the other stock's flow explains 4.5\% and 5.3\% of each stock's return variance beyond its own; trading the legs together
rather than in turn saves 19\% of the cost of the sale and 48\% of the pair trade, and an own-impact model misstates the cost by $-31\%$ and $+80\%$.
\textbf{18.} Both stocks together over the whole horizon, with a cost estimate that includes \omterm{def:mx:metaorders-latent-liquidity-and-cross-impact:cross}{cross-impact}.
\textbf{19.} Both legs simultaneously, so that each leg's \omterm{def:mx:metaorders-latent-liquidity-and-cross-impact:cross}{cross-impact} helps the other.
\textbf{20.} Its price moves with our trades in the stock we are selling, and so do our costs and our marks on it.
\end{solution}

\begin{solution}{iq:mx:metaorders-latent-liquidity-and-cross-impact:1}
Most liquidity is latent, and its density grows linearly away from the price. Moving the price by $I$ consumes $LI^2/2$ of it. So a metaorder of $Q$
moves the price by $\sqrt{2Q/L}$ when the book cannot refill in time.
\iqlookfor{Latent book; the integral; the fast regime.}
\end{solution}

\begin{solution}{iq:mx:metaorders-latent-liquidity-and-cross-impact:2}
The average price paid equals the price after completion; with square-root impact the price should settle near two thirds of the peak.
\iqlookfor{The condition; the two-thirds consequence.}
\end{solution}

\begin{solution}{iq:mx:metaorders-latent-liquidity-and-cross-impact:3}
Positions in correlated names, through \omterm{def:mx:metaorders-latent-liquidity-and-cross-impact:cross}{cross-impact} on their prices: marks and later sales suffer. Estimate \omterm{def:mx:metaorders-latent-liquidity-and-cross-impact:cross}{cross-impact} on my own trades or on \omterm{def:mx:the-limit-order-book:orders}{market
order} flow, by regressing returns of each name on the flows of the others.
\iqlookfor{\omterm{def:mx:metaorders-latent-liquidity-and-cross-impact:cross}{Cross-impact}; estimation on flows.}
\end{solution}

\begin{solution}{iq:mx:metaorders-latent-liquidity-and-cross-impact:4}
Regress interval returns on the vector of order-flow imbalances, with enough intervals and names grouped by sector; impose symmetry and positive
semi-definiteness (project the estimate), and consider an eigen-decomposition of the flows' covariance to reduce noise.
\iqlookfor{Regression on flows; symmetry and PSD; dimension reduction.}
\end{solution}

\begin{solution}{iq:mx:metaorders-latent-liquidity-and-cross-impact:5}
An explicit scheme needs $D\,dt/dx^2\le\tfrac12$ for stability; the boundary (the price) is found at each step as the zero of the density, interpolated
between cells; an implicit scheme removes the step limit at the cost of a linear solve.
\iqlookfor{The stability condition; boundary tracking.}
\end{solution}

\begin{solution}{iq:mx:metaorders-latent-liquidity-and-cross-impact:6}
With positive \omterm{def:mx:metaorders-latent-liquidity-and-cross-impact:cross}{cross-impact}, selling one leg lowers the other's price while buying it: the legs' cross terms offset. Trade them simultaneously and
balanced, and estimate the cost with the cross terms included.
\iqlookfor{Offsetting \omterm{def:mx:metaorders-latent-liquidity-and-cross-impact:cross}{cross-impact}; simultaneous execution.}
\end{solution}
